\chapter{Lehre}

Das Projekt ist auch ein Lerngegenstand. Dieses Kapitel beschreibt, wie es in Kursen
eingesetzt werden kann.

\section{Kurse}

\startitemize
\item {\bf DUA, Datenstrukturen und Algorithmen:} Hashing, Merkle-Graphen,
Myers-Diff und Graphsuche.
Der Kurs arbeitet mit dem Buch von Sedgewick und Wayne \cite[sedgewick2011algorithms];
\in{Tabelle}[tab:algorithms] ordnet die Algorithmen des Codes den Kapiteln des Buches zu.
\item {\bf AISE, AI Assisted Software Engineering:} Agenten-Workflows,
Branch-Regeln, Freigaben und Risikostufen.
\stopitemize

\section{Algorithmen im Code und im Buch}

Der Kurs DUA arbeitet mit dem Buch {\em Algorithms} von Sedgewick und Wayne
\cite[sedgewick2011algorithms].
\in{Tabelle}[tab:algorithms] nennt, welche Algorithmen und Datenstrukturen der Code von Stemma
heute einsetzt, wo man sie findet und in welchem Kapitel des Buches sie stehen.
Sie ist ein Wegweiser zum Nachlesen und zum Vergleichen mit der eigenen Lösung,
keine Beschreibung des Buches.

\startplacetable[title={Algorithmen im Code von Stemma},
  reference=tab:algorithms,location={here,force}]
\startxtable[frame=off,split=yes,option=stretch]
\startxtablehead
\startxrow[topframe=on,bottomframe=on]
\startxcell {\bf Algorithmus} \stopxcell
\startxcell {\bf Stelle im Code} \stopxcell
\startxcell {\bf Sedgewick und Wayne} \stopxcell
\stopxrow
\stopxtablehead
\startxtablenext
\startxrow[topframe=on,bottomframe=on]
\startxcell {\bf Algorithmus} \stopxcell
\startxcell {\bf Stelle im Code} \stopxcell
\startxcell {\bf Sedgewick und Wayne} \stopxcell
\stopxrow
\stopxtablenext
\startxrow
\startxcell Sortieren der Einträge eines Baums und der Branches eines Eintrags im Log \stopxcell
\startxcell \type{Tree::new} (\type{object.rs}), \type{Change::new} (\type{oplog.rs}) \stopxcell
\startxcell Kapitel 2, Sortieren \stopxcell
\stopxrow
\startxrow
\startxcell Binäre Suche in einer sortierten Tabelle von Bereichen \stopxcell
\startxcell \type{is_invisible_character} (\type{object.rs}) \stopxcell
\startxcell Abschnitt 1.1 und 3.1, binäre Suche \stopxcell
\stopxrow
\startxrow
\startxcell Geordnete Symboltabelle: Der Index hält die Pfade sortiert \stopxcell
\startxcell \type{Index} mit \type{BTreeMap} (\type{index.rs}) \stopxcell
\startxcell Abschnitt 3.1 bis 3.3, Symboltabellen und balancierte Suchbäume (ein B-Baum
verallgemeinert den 2-3-Baum) \stopxcell
\stopxrow
\startxrow
\startxcell Hashtabelle, um doppelte Pfade und doppelte Eltern zu erkennen \stopxcell
\startxcell \type{HashSet} in \type{replace_working_copy} (\type{snapshot.rs}) und
\type{Commit::new} \stopxcell
\startxcell Abschnitt 3.4, Hashtabellen \stopxcell
\stopxrow
\startxrow
\startxcell Tiefensuche durch Verzeichnisbäume, mit einem Stapel statt Rekursion, damit ein tiefer
Baum den Aufrufstapel nicht erschöpft \stopxcell
\startxcell \type{tracked_files} (\type{snapshot.rs}), \type{walk} (\type{track.rs}) \stopxcell
\startxcell Abschnitt 1.3, Stapel, und 4.1, Tiefensuche (ein Baum ist ein Graph ohne
Kreis) \stopxcell
\stopxrow
\startxrow
\startxcell Rückwärtiges Durchlaufen der Einträge, um das jüngste Vorkommen je
Zeiger zu finden \stopxcell
\startxcell \type{outcomes} (\type{oplog.rs}) \stopxcell
\startxcell Abschnitt 1.3 und 3.1, lineares Durchlaufen und Symboltabelle \stopxcell
\stopxrow
\stopxtable
\stopplacetable

Drei Dinge aus den Kursthemen sind noch nicht im Code oder stehen nicht im Buch:

\startitemize
\item {\bf Kryptografische Hashfunktion:} Die Objekt-IDs sind BLAKE3-Hashes des Inhalts
(ADR-001).
Das Buch behandelt Hashtabellen und ihre Hashfunktionen, nicht kryptografische Hashfunktionen.
Für die IDs gilt eine andere Anforderung: Man darf zwei verschiedene Inhalte
nicht mit gleichem Hash finden können.
\item {\bf Merkle-Graph:} Bäume und Commits nennen ihre Teile über die ID des Inhalts
(\type{object.rs}).
Die Struktur ist ein gerichteter Graph ohne Kreis (Abschnitt 4.2 des Buches);
den Merkle-Baum als Begriff kennt das Buch nicht.
Den Commit-Graphen durchläuft erst Phase 3, für die Suche nach dem gemeinsamen Vorfahren
(Abschnitt 4.1 und 4.2, Breitensuche und Tiefensuche).
\item {\bf Myers-Diff:} Er kommt in Phase 2 und steht nicht im Buch.
Am nächsten liegt das Problem der längsten gemeinsamen Teilfolge,
das sich mit dynamischer Programmierung lösen lässt;
Myers löst es in einer Zeit, die von der Zahl der Unterschiede abhängt.
\stopitemize

\section{Übungen je Phase}

Jede Phase lässt sich als Übung formulieren. Studierende implementieren einen Teil
selbst und vergleichen ihn mit einer Referenzimplementierung. Die Phase 4 eignet
sich besonders, weil Freigaben und Selbstfreigabeverbot eine Prüfung durch Agenten
und Menschen verlangen.

\section{Offene Fragen}

\startitemize
\item Lizenz des Lehrmaterials und Umgang mit fremden Quellen.
\item Bewertungskriterien für Code, der mit Agenten entstanden ist.
\item Trennung von Studierendennamespaces und Dozierendenrolle.
\stopitemize
